% =========================================================
% English-to-Logic Translation and Scope Ambiguity
% =========================================================

\subsection*{English to Logic and Scope Ambiguity}

% BEGIN GENERATED ARTIFACT: def:translation-key-predicate-logic
\begin{definitionbox}{Definition (Translation Key)}
\begin{definition}[Translation Key]
\label{def:translation-key-predicate-logic}
A \emph{translation key} is a finite record that fixes a domain of discourse and
assigns intended English readings to the constant symbols, predicate symbols,
and relation symbols used in a translation.
\end{definition}
\end{definitionbox}

\begin{remark*}[Standard quantified statement]
\[
\begin{aligned}
&K \text{ is a translation key}\\
&\Longleftrightarrow
K \text{ specifies a domain and readings for each non-logical symbol used.}
\end{aligned}
\]
\end{remark*}

\begin{remark*}[Predicate reading]
\(\operatorname{TranslationKey}(K)\) means that \(K\) supplies the domain and
symbol readings needed to parse a translated formula.
\end{remark*}

\begin{remark*}[Interpretation]
A translation key prevents hidden changes of meaning. The same symbol must keep
the same reading throughout one exercise unless the key is explicitly replaced.
\end{remark*}

\begin{remark*}[Examples]
For the English sentence ``Every applicant submitted some form,'' a key may
take the domain to be people and forms, with \(A(x)\) for ``\(x\) is an
applicant,'' \(F(y)\) for ``\(y\) is a form,'' and \(S(x,y)\) for ``\(x\)
submitted \(y\).''
\end{remark*}

\begin{remark*}[Non-Examples]
Using \(P(x)\) once for ``\(x\) is a person'' and later for ``\(x\) passed'' is
not a translation key; it is a change of vocabulary inside one translation.
\end{remark*}

\begin{dependencies}
\begin{itemize}
  \item \hyperref[def:constant-symbol]{Constant Symbol}
  \item \hyperref[def:predicate-symbol]{Predicate Symbol}
\end{itemize}
\end{dependencies}
% END GENERATED ARTIFACT: def:translation-key-predicate-logic

% BEGIN GENERATED ARTIFACT: def:domain-of-discourse-predicate-logic
\begin{definitionbox}{Definition (Domain of Discourse)}
\begin{definition}[Domain of Discourse]
\label{def:domain-of-discourse-predicate-logic}
The \emph{domain of discourse} for a translation is the collection of objects
over which the variables of the translated formulas range.
\end{definition}
\end{definitionbox}

\begin{remark*}[Standard quantified statement]
\[
D \text{ is the domain of discourse }
\Longleftrightarrow
\text{ each variable in the translation ranges over objects in }D.
\]
\end{remark*}

\begin{remark*}[Predicate reading]
\(\operatorname{DomainOfDiscourse}(D,K)\) means that \(D\) is the domain fixed
by the translation key \(K\).
\end{remark*}

\begin{remark*}[Interpretation]
The domain decides what the quantifiers are allowed to talk about. If the
domain is people, then an exam must either be named by a separate sort-like
predicate or represented by a different translation choice.
\end{remark*}

\begin{remark*}[Examples]
For ``All cats are mammals,'' one may use the domain of all animals and write
\(\forall x\,(C(x)\to M(x))\). With the smaller domain of cats, the same English
claim may be written \(\forall x\,M(x)\).
\end{remark*}

\begin{remark*}[Non-Examples]
Writing \(\forall x\,M(x)\) while silently changing the domain from animals to
cats halfway through the exercise is not a valid translation.
\end{remark*}

\begin{dependencies}
\begin{itemize}
  \item \hyperref[def:translation-key-predicate-logic]{Translation Key}
  \item \hyperref[def:variable]{Variable}
\end{itemize}
\end{dependencies}
% END GENERATED ARTIFACT: def:domain-of-discourse-predicate-logic

% BEGIN GENERATED ARTIFACT: def:translation-predicate-relation-symbol
\begin{definitionbox}{Definition (Predicate and Relation Symbols in a Translation)}
\begin{definition}[Predicate and Relation Symbols in a Translation]
\label{def:translation-predicate-relation-symbol}
In a translation key, a unary predicate symbol names a property of one object,
and an \(n\)-ary relation symbol names a relation among \(n\) objects.
\end{definition}
\end{definitionbox}

\begin{remark*}[Standard quantified statement]
\[
\begin{aligned}
&P \text{ is unary in }K &&\Longleftrightarrow P(x) \text{ has one argument,}\\
&R \text{ is }n\text{-ary in }K &&\Longleftrightarrow R(x_1,\ldots,x_n)
  \text{ has }n\text{ arguments.}
\end{aligned}
\]
\end{remark*}

\begin{remark*}[Predicate reading]
\(\operatorname{TranslationPredicateSymbol}(P,K)\) records that \(P\) has a
fixed predicate reading in \(K\), and \(\operatorname{TranslationRelationSymbol}(R,n,K)\)
records that \(R\) is used with arity \(n\) in \(K\).
\end{remark*}

\begin{remark*}[Interpretation]
Object-language predicate symbols such as \(P,Q,R\) are symbols inside the
formal language. Metalinguistic predicate readings such as
\(\operatorname{Term}(t)\) or \(\operatorname{AtomicFormula}(\varphi)\) are
repository descriptions of syntax; they are not formulas being translated.
\end{remark*}

\begin{remark*}[Examples]
\(A(x)\) may mean ``\(x\) is an athlete.'' \(L(x,y)\) may mean ``\(x\) likes
\(y\).'' The arity is part of the key: \(L(x)\) and \(L(x,y)\) are not the same
symbol use.
\end{remark*}

\begin{remark*}[Non-Examples]
Writing \(L(x,y,z)\) after declaring \(L(x,y)\) as ``\(x\) likes \(y\)'' is an
arity error.
\end{remark*}

\begin{dependencies}
\begin{itemize}
  \item \hyperref[def:predicate-symbol]{Predicate Symbol}
  \item \hyperref[def:atomic-formula]{Atomic Formula}
\end{itemize}
\end{dependencies}
% END GENERATED ARTIFACT: def:translation-predicate-relation-symbol

% BEGIN GENERATED ARTIFACT: def:translation-constants-variables
\begin{definitionbox}{Definition (Constants and Variables in a Translation)}
\begin{definition}[Constants and Variables in a Translation]
\label{def:translation-constants-variables}
A \emph{constant} names a fixed object from the translation key, while a
\emph{variable} marks a place bound by a quantifier or left open for later
binding.
\end{definition}
\end{definitionbox}

\begin{remark*}[Standard quantified statement]
\[
\begin{aligned}
&c \text{ is a translation constant} &&\Longleftrightarrow c
  \text{ names one fixed object,}\\
&x \text{ is a translation variable} &&\Longleftrightarrow x
  \text{ is available for quantification or free occurrence.}
\end{aligned}
\]
\end{remark*}

\begin{remark*}[Predicate reading]
\(\operatorname{TranslationConstant}(c,K)\) means that \(c\) names a fixed
object in \(K\), and \(\operatorname{TranslationVariable}(x,K)\) means that
\(x\) is used as a variable in translations governed by \(K\).
\end{remark*}

\begin{remark*}[Interpretation]
Constants do not range. Variables do range when bound by quantifiers. Confusing
the two usually leaves a formula with the wrong dependency pattern.
\end{remark*}

\begin{remark*}[Examples]
If \(a\) names Alice, then \(P(a)\) says Alice is a participant. The formula
\(\forall x\,(S(x)\to P(x))\) uses \(x\) as a bound variable.
\end{remark*}

\begin{remark*}[Non-Examples]
The expression \(\forall a\,S(a)\) is not appropriate if \(a\) has already been
declared as a constant naming Alice.
\end{remark*}

\begin{dependencies}
\begin{itemize}
  \item \hyperref[def:constant-symbol]{Constant Symbol}
  \item \hyperref[def:variable]{Variable}
\end{itemize}
\end{dependencies}
% END GENERATED ARTIFACT: def:translation-constants-variables

% BEGIN GENERATED ARTIFACT: prop:universal-translation-pattern
\begin{propositionbox}{Proposition (Universal Translation Pattern)}
\begin{proposition}[Universal Translation Pattern]
\label{prop:universal-translation-pattern}
In a broad domain, an English statement of the form ``All \(A\) are \(B\)''
is translated by
\[
\forall x\,(A(x)\to B(x)).
\]

\noindent\hyperref[prf:universal-translation-pattern]{Go to proof.}
\end{proposition}
\end{propositionbox}

\begin{remark*}[Standard quantified statement]
\[
\forall x\,(A(x)\to B(x)).
\]
\end{remark*}

\begin{remark*}[Predicate reading]
\(\operatorname{UniversalTranslationPattern}(A,B)\) records the implication
schema for translating universal restricted claims.
\end{remark*}

\begin{remark*}[Interpretation]
The antecedent \(A(x)\) restricts attention to the relevant objects. Objects
outside \(A\) do not have to satisfy \(B\).
\end{remark*}

\begin{remark*}[Exposition]
``All cats are mammals'' becomes \(\forall x\,(C(x)\to M(x))\). ``No cats are
reptiles'' becomes \(\forall x\,(C(x)\to \neg R(x))\). The form
\(\forall x\,(C(x)\land M(x))\) says every object in the domain is both a cat
and a mammal, which is far stronger than the intended restricted claim. In a
broad domain, ``all,'' ``every,'' and ``any'' usually translate restricted
claims with implication.
\end{remark*}

\begin{dependencies}
\begin{itemize}
  \item \hyperref[def:universal-q]{Universal Formula}
  \item \hyperref[def:molecular]{Molecular Formula}
\end{itemize}
\end{dependencies}
% END GENERATED ARTIFACT: prop:universal-translation-pattern

% BEGIN GENERATED ARTIFACT: prop:existential-translation-pattern
\begin{propositionbox}{Proposition (Existential Translation Pattern)}
\begin{proposition}[Existential Translation Pattern]
\label{prop:existential-translation-pattern}
In a broad domain, an English statement of the form ``Some \(A\) are \(B\)''
is translated by
\[
\exists x\,(A(x)\land B(x)).
\]

\noindent\hyperref[prf:existential-translation-pattern]{Go to proof.}
\end{proposition}
\end{propositionbox}

\begin{remark*}[Standard quantified statement]
\[
\exists x\,(A(x)\land B(x)).
\]
\end{remark*}

\begin{remark*}[Predicate reading]
\(\operatorname{ExistentialTranslationPattern}(A,B)\) records the conjunction
schema for translating existential restricted claims.
\end{remark*}

\begin{remark*}[Interpretation]
The witness must satisfy both restrictions at once. Existence is not merely a
conditional promise about what would happen if an \(A\) existed.
\end{remark*}

\begin{remark*}[Exposition]
``Some cats are black'' becomes \(\exists x\,(C(x)\land B(x))\). ``Some cats are
not black'' becomes \(\exists x\,(C(x)\land \neg B(x))\). The form
\(\exists x\,(C(x)\to B(x))\) may be true because the witness is not a cat, so
it does not translate ``Some cats are black.'' In a broad domain, ``some,''
``there is,'' and ``at least one'' usually translate restricted claims with
conjunction.
\end{remark*}

\begin{dependencies}
\begin{itemize}
  \item \hyperref[def:existential-q]{Existential Formula}
  \item \hyperref[def:molecular]{Molecular Formula}
\end{itemize}
\end{dependencies}
% END GENERATED ARTIFACT: prop:existential-translation-pattern

% BEGIN GENERATED ARTIFACT: prop:conditional-translation-pattern
\begin{propositionbox}{Proposition (Conditional Translation Pattern)}
\begin{proposition}[Conditional Translation Pattern]
\label{prop:conditional-translation-pattern}
An English statement of the form ``If \(A\), then \(B\)'' is translated by
\[
A\to B,
\]
with quantifiers placed according to the variables occurring in \(A\) and \(B\).

\noindent\hyperref[prf:conditional-translation-pattern]{Go to proof.}
\end{proposition}
\end{propositionbox}

\begin{remark*}[Standard quantified statement]
\[
A\to B.
\]
\end{remark*}

\begin{remark*}[Predicate reading]
\(\operatorname{ConditionalTranslationPattern}(A,B)\) records the implication
schema for conditional English.
\end{remark*}

\begin{remark*}[Interpretation]
Conditionals are especially common inside universal translations. The
conditional separates the restricting condition from the asserted consequence.
\end{remark*}

\begin{remark*}[Exposition]
``Every applicant who is eligible is interviewed'' becomes
\(\forall x\,((A(x)\land E(x))\to I(x))\). The form
\(\forall x\,(A(x)\land E(x)\land I(x))\) says every object is an eligible
applicant who is interviewed, which is not the intended conditional claim.
\end{remark*}

\begin{dependencies}
\begin{itemize}
  \item \hyperref[prop:universal-translation-pattern]{Universal Translation Pattern}
\end{itemize}
\end{dependencies}
% END GENERATED ARTIFACT: prop:conditional-translation-pattern

% BEGIN GENERATED ARTIFACT: prop:negation-translation-pattern
\begin{propositionbox}{Proposition (Negation Translation Pattern)}
\begin{proposition}[Negation Translation Pattern]
\label{prop:negation-translation-pattern}
The negation of a translated quantified statement is obtained by negating the
whole formula and then pushing the negation inward using the quantifier
negation laws.

\noindent\hyperref[prf:negation-translation-pattern]{Go to proof.}
\end{proposition}
\end{propositionbox}

\begin{remark*}[Standard quantified statement]
\[
\begin{aligned}
\neg\forall x\,\varphi &\Longleftrightarrow \exists x\,\neg\varphi,\\
\neg\exists x\,\varphi &\Longleftrightarrow \forall x\,\neg\varphi.
\end{aligned}
\]
\end{remark*}

\begin{remark*}[Predicate reading]
\(\operatorname{NegationTranslationPattern}(\varphi,\psi)\) records that
\(\psi\) is obtained from \(\neg\varphi\) by pushing negation through the
quantifiers and connectives.
\end{remark*}

\begin{remark*}[Interpretation]
The negation of ``all'' is ``some not.'' The negation of ``some'' is ``none.''
For restricted universal claims, negating the implication produces a witness
that satisfies the restriction and fails the consequent.
\end{remark*}

\begin{remark*}[Exposition]
The negation of \(\forall x\,(A(x)\to B(x))\) is
\(\exists x\,(A(x)\land \neg B(x))\). The negation of
\(\exists x\,(A(x)\land B(x))\) is \(\forall x\,(A(x)\to \neg B(x))\). The
negation of \(\forall x\,(A(x)\to B(x))\) is not
\(\forall x\,(A(x)\to \neg B(x))\); that latter formula says every \(A\) fails
to be \(B\), not merely that some counterexample exists.
\end{remark*}

\begin{dependencies}
\begin{itemize}
  \item \hyperref[thm:qneg]{Quantifier Negation Laws}
  \item \hyperref[prop:universal-translation-pattern]{Universal Translation Pattern}
  \item \hyperref[prop:existential-translation-pattern]{Existential Translation Pattern}
\end{itemize}
\end{dependencies}
% END GENERATED ARTIFACT: prop:negation-translation-pattern

% BEGIN GENERATED ARTIFACT: def:scope-ambiguity-translation
\begin{definitionbox}{Definition (Scope Ambiguity)}
\begin{definition}[Scope Ambiguity]
\label{def:scope-ambiguity-translation}
A sentence has \emph{scope ambiguity} when the same ordinary-language wording
admits two or more translations with different quantifier nesting or different
connective scope.
\end{definition}
\end{definitionbox}

\begin{remark*}[Standard quantified statement]
\[
\begin{aligned}
&E \text{ has scope ambiguity}\\
&\Longleftrightarrow
\text{ there exist distinct candidate translations }\varphi,\psi
\text{ of }E\text{ with different scope.}
\end{aligned}
\]
\end{remark*}

\begin{remark*}[Predicate reading]
\(\operatorname{ScopeAmbiguity}(E,\varphi,\psi)\) means that the English
sentence \(E\) can reasonably be translated by \(\varphi\) or by \(\psi\), and
the difference is a scope difference.
\end{remark*}

\begin{remark*}[Interpretation]
Scope is not cosmetic. Changing the outer quantifier can change whether a
single witness works for everyone or whether each object may have its own
witness.
\end{remark*}

\begin{remark*}[Exposition]
The formulas \(\forall x\exists y\,R(x,y)\) and
\(\exists y\forall x\,R(x,y)\) are not interchangeable. The first allows the
witness \(y\) to depend on \(x\); the second requires one \(y\) that works for
every \(x\). Symbols such as \(A(x)\), \(B(x)\), and \(R(x,y)\) are
object-language formula patterns used in translations. Expressions such as
\(\operatorname{TranslationKey}(K)\), \(\operatorname{Term}(t)\), and
\(\operatorname{AtomicFormula}(\varphi)\) are metalinguistic predicate readings
used by the project to describe mathematical objects.
\[
\begin{array}{lll}
\text{All }A\text{ are }B
  &\leadsto& \forall x\,(A(x)\to B(x))\\
\text{Some }A\text{ are }B
  &\leadsto& \exists x\,(A(x)\land B(x))\\
\text{No }A\text{ are }B
  &\leadsto& \forall x\,(A(x)\to \neg B(x))\\
\text{Some }A\text{ are not }B
  &\leadsto& \exists x\,(A(x)\land \neg B(x))\\
\text{Every }A\text{ has a }B
  &\leadsto& \forall x\,(A(x)\to \exists y\,(B(y)\land R(x,y)))\\
\text{There is a }B\text{ for every }A
  &\leadsto& \exists y\,(B(y)\land \forall x\,(A(x)\to R(x,y))).
\end{array}
\]
\end{remark*}

\begin{remark*}[Examples]
``Every \(A\) has a \(B\)'' is usually
\(\forall x\,(A(x)\to \exists y\,(B(y)\land R(x,y)))\). ``There is a \(B\) for
every \(A\)'' is usually
\(\exists y\,(B(y)\land \forall x\,(A(x)\to R(x,y)))\). The bare patterns
\(\forall x\exists y\,R(x,y)\) and \(\exists y\forall x\,R(x,y)\) make the same
scope contrast visible.
\end{remark*}

\begin{remark*}[Non-Examples]
Renaming a bound variable from \(x\) to \(z\) does not create a scope ambiguity.
The scope is unchanged.
\end{remark*}

\begin{dependencies}
\begin{itemize}
  \item \hyperref[def:quantifier-scope]{Quantifier Scope}
  \item \hyperref[def:q-alt]{Quantifier Order}
\end{itemize}
\end{dependencies}
% END GENERATED ARTIFACT: def:scope-ambiguity-translation
